% !TeX program = pdflatex
% ALIUS Bulletin standalone facsimile source.
% Compile from the ALIUS-bulletin project root for exact visual fidelity.
% If the shared PDF asset is absent, this file falls back to extracted text.
\providecommand{\ALIUSRootPrefix}{}

\ifdefined\ALIUSIssueBuild
  \IfFileExists{\ALIUSRootPrefix Shared-assets/original-pdfs/Issue02/alius-bulletin-issue-02-fox-wandering-along-the-spectrum-of-spontaneous-thinking-dreaming-meditation.pdf}{%
    \includepdf[pages=-,fitpaper=true]{\ALIUSRootPrefix Shared-assets/original-pdfs/Issue02/alius-bulletin-issue-02-fox-wandering-along-the-spectrum-of-spontaneous-thinking-dreaming-meditation.pdf}%
  }{%
    \typeout{ALIUS warning: missing source PDF Shared-assets/original-pdfs/Issue02/alius-bulletin-issue-02-fox-wandering-along-the-spectrum-of-spontaneous-thinking-dreaming-meditation.pdf}%
  }%
  \expandafter\endinput
\fi

\documentclass[a4paper,11pt]{article}
\usepackage[margin=22mm]{geometry}
\usepackage{pdfpages}
\usepackage[hidelinks]{hyperref}
\usepackage{microtype}
\usepackage{parskip}
\title{Wandering along the spectrum of spontaneous thinking: Dreaming, meditation, mind-wandering, and well-being}
\author{Kieran Fox Interviewed by Matthieu Koroma}
\date{ALIUS Bulletin 2}
\begin{document}
\IfFileExists{\ALIUSRootPrefix Shared-assets/original-pdfs/Issue02/alius-bulletin-issue-02-fox-wandering-along-the-spectrum-of-spontaneous-thinking-dreaming-meditation.pdf}{%
  \includepdf[pages=-,fitpaper=true]{\ALIUSRootPrefix Shared-assets/original-pdfs/Issue02/alius-bulletin-issue-02-fox-wandering-along-the-spectrum-of-spontaneous-thinking-dreaming-meditation.pdf}%
}{%
  \maketitle
  \section*{Metadata}
  \textbf{Piece type:} interview\par
  \textbf{DOI:} 10.34700/m6kp-c494\par
  \textbf{Keywords:} spontaneous thought, mind-wandering, default network, mindfulness meditation, current concerns\par
  \section*{Abstract}
  This interview with Kieran Fox examines spontaneous thinking as an umbrella covering mind-wandering, dreaming, and meditation, and its connections to creativity and well-being. Fox recounts how his work with Kalina Christoff led him to map the parallels between mind-wandering and dreaming. The conversation lays out the spontaneous thought framework, in which cognition sits on a spectrum between totally constrained thought and disorganized thinking approaching acute psychosis, with mind-wandering, dreaming, and creative thinking falling between these extremes. Fox argues that most mind-wandering proceeds below meta-awareness and compares the stream of inner thought to other perceptual channels filtered by attention. He treats mind-wandering and dreaming as matters of degree, with sensory decoupling and increased activity in default, memory, and visual networks accounting for the more immersive character of dream experience. Further sections examine the determinants of thought content, including dominant visual, auditory, and somatosensory modalities, mildly positive average emotional valence, and Eric Klinger's notion of current concerns. The conversation turns to mindfulness meditation as a way to investigate spontaneous thought without disturbing it, citing work on long-term practitioners that locates the medial temporal lobe and default network among the antecedent activations of arising thoughts. Fox then considers the adaptive value of mind-wandering, flow states, and creativity as a special case of focus on current concerns, before treating post-traumatic stress disorder and depressive rumination as distortions of spontaneous thinking. The interview closes on neuroimaging evidence linking mind-wandering to the default network and on future work into the stability and malleability of thought patterns.\par
  \section*{Extracted Text}
  \begingroup
  \small
  \noindent K. Fox - Wandering along the spectrum of spontaneous thinking ALIUS Bulletin n2 (2018) aliusresearch.org/bulletin 1 Wandering along the spectrum of spontaneous thinking Dreaming, meditation, mind-wandering, and well-being An interview with Kieran Fox by Matthieu Koroma Kieran Fox kcrfox@stanford.edu Department of Neurology and Neurological Sciences Stanford University, USA Matthieu Koroma mkoroma@ens-paris-saclay.fr LSCP, Departement d'etudes cognitives ENS/EHESS/PSL Research University, Paris, France Citation: Fox, K. \& Koroma, M. (2018). Wandering along the spectrum of spontaneous thinking: dreaming, meditation, mind-wandering, and well-being. An interview with Kieran Fox. ALIUS Bulletin, 2, 1-15. What has triggered your scientific interest in the study of mind-wandering and other associated phenomena? I really owe my interest in mind-wandering specifically and spontaneous thought more generally to my PhD supervisor, Kalina Christoff. When I started my PhD with Kalina, my interest was in studying the brain basis of meditation and neurofeedback. I thought of mind-wandering the way many other people did at the time: as an annoyance, something that gets in the way of meditation and distracts you from more valuable work. I already had a longstanding interest in sleep and dreaming, but at the time I didn't see all the parallels (both neural and psychological) between dreaming and mind-wandering; exploring those parallels was a big part of my graduate work (Domhoff \& Fox, 2015; Fox \& Christoff, 2014; Fox, Nijeboer, Solomonova, Domhoff, \& Christoff, 2013). Talking with Kalina and reading her prior research showed me how interesting and important mind-wandering really is: how it's related to creativity, dreaming, personality, even mental illness. Our mutual interest in this topic recently led us to edit a book-length treatment of the subject which is due out soon (May 2018 : http://amzn.to/2rLQW9d). Mind-wandering has been the focus of recent research with a diversity of approaches and definitions. It has been proposed that the notion of mind-wandering could be considered as an umbrella term for a range of phenomena with family- resemblances studied using different approaches, rather than being defined as a K. Fox - Wandering along the spectrum of spontaneous thinking ALIUS Bulletin n2 (2018) aliusresearch.org/bulletin 2 cognitive phenomenon per se. What is the definition or approach of mind- wandering that you endorse in your work? Our main effort at conceptualizing these states uses spontaneous thought as the umbrella term (Christoff, Irving, Fox, Spreng, \& Andrews-Hanna, 2016). By "spontaneous" we mean cognition that is relatively unconstrained-either by bottom-up constraints like affective or perceptual salience, or top-down constraints like the executive control of attention. In this framework, mind-wandering falls somewhere in the center of a spectrum. On one end of this spectrum you have totally unconstrained thought; no one really knows what this would look like, but the closest we can imagine is acute psychosis-there is thought and perception and mental content, but it is totally disorganized and nothing is connected, no thought is really related to any other. At the other extreme, you have totally constrained thought, for instance if you were totally focused on writing a difficult exam. Cognition like dreaming, mind-wandering, and creative thinking instead fall in the middle-they are more spontaneous forms of thought than, say, writing a demanding exam, but they are still far more constrained and coherent than, say, psychosis. This is the current thinking of my colleagues and myself, but, of course, these states are difficult to study and define and we hope to make more progress in the future (see Figure 1, reproduced with permission from Christoff, Irving, Fox, Spreng, \& Andrews-Hanna, 2016). Is mind-wandering necessarily a conscious phenomenon or can there be a subconscious form of mind-wandering, with a drifting in the activation of different cognitive networks without being consciously accessed? My own view, and I think this is amply supported by the empirical evidence, is that mind-wandering can definitely be unconscious-in fact, I think it's safe to say the majority of mind-wandering is below the level of full awareness. For example, in laboratory studies of mind-wandering, when randomly-timed "thought probes" catch people mind-wandering, we can also ask whether people were conscious or "meta-aware" of their mind-wandering. In these studies, meta-awareness of mind- wandering is only reported about half of the time (Christoff, Gordon, Smallwood, Smith, \& Schooler, 2009; Fox \& Christoff, 2015; Schooler et al., 2011). And this is even during studies where people know they are going to be asked about mind- wandering; even then, they only seem to notice about 50\%. Assuming you are By 'spontaneous' we mean cognition that is relatively unconstrained-either by bottom-up constraints like affective or perceptual salience, or top-down constraints like the executive control of attention. " " K. Fox - Wandering along the spectrum of spontaneous thinking ALIUS Bulletin n2 (2018) aliusresearch.org/bulletin 3 K. Fox - Wandering along the spectrum of spontaneous thinking ALIUS Bulletin n2 (2018) aliusresearch.org/bulletin 4 persuaded by the parallels we have drawn between mind-wandering when awake and dreaming while asleep (Domhoff \& Fox, 2015; Fox et al., 2013), then dreaming offers an even more compelling case: if you come to a sleep laboratory and are awakened when your brain is in REM sleep, there is an extremely high chance (around 80\% or higher) that you will report a detailed, immersive dream experience. We go through four to six REM periods each night, and there is certainly dreaming taking place in other sleep stages as well (Fox \& Girn, in press; Nielsen, 2000). Even using a very conservative estimate, we are probably having 10 dream experiences each night at the very least, yet the average person remembers none of this whatsoever, and even people very interested in recording and analyzing their dreams find them difficult to recall. So I think the evidence that detailed and complex spontaneous thought processes can take place without any concurrent awareness or subsequent recall is overwhelming, and in fact this raises some very deep questions about why these phenomena take place at all-why are there rich psychological/subjective correlates of these brain processes that we so seldom notice and whose function (if any) is very difficult to determine? If you consider mind-wandering as a conscious phenomenon, would you consider mind-wandering as a proper conscious state (in a sense that mind- wandering characterizes a way of being conscious (Bayne et al., 2016), like for example a meditative state)? I don't think of mind-wandering as a conscious state-I think of these processes as more or less ongoing, below the level of awareness, competing with other inputs and signals in the brain for our attention. We can tune in and pay attention to them, or not, and sometimes the thoughts will be strong enough or emotionally salient enough to grab our attention even when we don't want them to. I think of the stream of inner thought in a way similar to other perceptual channels; for instance, you are constantly receiving a stream of auditory information, even when you're asleep, but your brain is very good at blocking out probably 99\% of this information as totally irrelevant, and you never become conscious of it. But this doesn't mean that your ears are not receiving the sounds, that the sound is not being transduced and conducted along the auditory nerve, and processed at least at some low level in the brain. I suspect the brain is constantly generating thoughts, imagery, and so on at a "subthreshold" level as well, and noticing it is more a matter of this content catching I think the evidence that detailed and complex spontaneous thought processes can take place without any concurrent awareness or subsequent recall is overwhelming. " " K. Fox - Wandering along the spectrum of spontaneous thinking ALIUS Bulletin n2 (2018) aliusresearch.org/bulletin 5 our attention and becoming illuminated by our conscious awareness than of entering a particular conscious state where mind-wandering then starts or is allowed to take place. According to work from you and your colleagues (Fox et al, 2013, Domhoff \& Fox, 2015, Christoff et al., 2016, Fox \& Girn, in press, Domhoff, 2018), the difference between mind-wandering and dreaming can be seen as a matter of degree, dreams being a more intense form of mind-wandering. This intensification of internally- generated cognition can be temptingly explained by the specific situation of sleep, in which the brain is disconnected from its environment. As such, the disconnection of the dreamer can be seen as a condition for the occurrence and intensity of mind- wandering through the restriction on the relay of sensory information. Conversely, disconnection of the dreamer can be seen as a consequence of the increase in internal activity that competes for attentional resources with externally-oriented networks (Nir \& Tononi, 2010). How is regulating the processing of external information occurring during episodes of spontaneous thinking? I think there is evidence that both explanations are correct. On the one hand, there is a clear dampening and near-blockade of many sensory inputs during sleep, suggesting that internally-generated channels of information have less "competition" for access to conscious awareness. And on the other hand, although brain metabolism tends to decrease in the NREM sleep stages, PET research suggests that the brain's energy usage equals (Braun et al., 1997; Madsen et al., 1991; Maquet et al., 1990) or perhaps even exceeds (Buchsbaum et al., 1989; Heiss, Pawlik, Herholz, Wagner, \& Wienhard, 1985) that of waking rest during REM sleep, when we know dreaming is most likely to occur. My colleagues and I have suggested that this has specific effects at both the psychological and neural level. In the brain, the increased activation appears to be preferentially localized to the default, memory, and visual networks (Fox et al., 2013), and we think this can help account for the psychological differences, namely an intensification of visual imagery, lengthy narratives, and so on. In waking spontaneous thought, there is a large body of work, especially from EEG, that mind-wandering is associated with decreased attention to the external environment (Kam, Dao, Stanciulescu, Tildesley, \& Handy, 2013; Kam \& Handy, 2013). So, to my mind the mechanisms appear quite similar: short, brief disconnections from external sensory inputs while awake can lead to correspondingly brief and relatively mild spontaneous thoughts; and similarly, much more drastic decoupling from the external world during sleep allows for much longer, more intense, and more immersive forms of spontaneous thought to take place. In a recent article that you co-authored, mind-wandering is characterized within the spontaneous thinking framework as a thought process in which content is weakly constrained by automatic processes and task-related cognitive control (Christoff et K. Fox - Wandering along the spectrum of spontaneous thinking ALIUS Bulletin n2 (2018) aliusresearch.org/bulletin 6 al, 2016). What, then, are the determinants of thought contents during mind- wandering? That is a very complicated question and the short answer is that we still know very little about what drives and determines the content of mind-wandering. But so- called "thought sampling" or "experience sampling" studies have allowed us to draw some broad conclusions. For instance, we know about the modalities in which mind- wandering tends to occur: these thoughts are very likely to be visual (in the form of imagery), auditory (as in "talking to yourself" or imagining conversations with other people), and somatosensory (thoughts about how your body feels) (reviewed in Fox, Andrews-Hanna, \& Christoff, 2016; Fox et al., 2013) These trends appear to hold across different populations in different countries, suggesting that this is a culture- independent neurophysiological process: human brains, in general, don't tend to think in smells or tastes; rather, thoughts take the form of visual imagery and imagined speech (most of the time, anyway). Another very robust finding is that mind-wandering-like thought contains a lot of emotional material; the majority of thoughts have some affective component, and on average they tend to be mildly positive. Contrary to widespread popular opinion, people overall are thinking about things they feel to be neutral or pleasant/positive (Fox et al., under review; Fox, Thompson, Andrews-Hanna, \& Christoff, 2014). But probably the most important determinant is what we call "current concerns", meaning what the individual cares about most at any given time, be it major things like a job interview or a sick relative, or more trivial concerns like what to get at the grocery store for tonight's dinner. The evidence is overwhelming that people think about their own personal goals and concerns a huge amount of the time (Klinger, 2008, 2013; Klinger \& Cox, 2004, 2011), which suggests that mind-wandering is not nearly as random or pointless as it is often made out to be. Instead, it seems as if there is a very clear functional role, that the brain is frequently working on processing and tackling those things that most concern us, even if we don't yet understand how this "work" is taking place. But I don't think it's a coincidence that insights into, and solutions to, our problems also usually come to us just as spontaneously and unexpectedly. Mindfulness meditation is a practice based on the careful observation of the train of thoughts through the exercise of meta-awareness. Through adopting "an open, non- judgmental metacognitive stance" (Fox \& Christoff, 2014), spontaneous mental activity can be accessed without trying to react to or control the content of the experience. Mindfulness meditation seems thus to offer a privileged situation for introspecting on mind-wandering. Could you elaborate on how such practice can offer insights into the nature and dynamics of spontaneous thinking? A major issue that many of us worry about in mind-wandering research is, "Are we changing the content and dynamics of spontaneous thought by asking people to K. Fox - Wandering along the spectrum of spontaneous thinking ALIUS Bulletin n2 (2018) aliusresearch.org/bulletin 7 observe and report on it?" The answer must almost certainly be "Yes." William James, in his Principles of Psychology, offered an apt simile for this problem, saying that trying to stop the mind at work to observe its functioning was like trying to stop a spinning top to more carefully investigate what its motion is like. By interrupting and analyzing the process, you alter it and lose something essential. This is where meditation broadly and mindfulness more specifically can potentially be helpful. Long-term mindfulness practitioners spend thousands of hours trying to observe their thoughts dispassionately, without reacting or altering things in any way. Whether they are successful or not is an open question and I think a very difficult one to answer, but there is some tentative evidence that long-term practitioners show more accurate and unbiased introspection (Fox et al., 2012; Sze et al., 2010). In a recent study I was involved in (Ellamil et al., 2016), we tried to harness these heightened introspective abilities in long-term practitioners to see if they could tell us about the exact moment when a spontaneous thought was arising in their minds. By using this timestamp as a marker, we could investigate what the brain was doing just prior to the conscious awareness of a thought, and try to infer how the brain was generating spontaneous thoughts in the first place. We found that the medial temporal lobe and default network regions were most prominent among the antecedent neural activations, suggesting that these areas play a key generative role. Other brain areas we know to be involved in spontaneous thought, such as prefrontal executive areas, came online a couple of seconds later, suggesting that they play a different role (perhaps in guiding or selecting how these thoughts are interpreted and responded to). This is just one example of how contemplative practitioners can help us answer subtle and tricky questions about spontaneous thought specifically, and human cognition more generally. I think this is a potentially very fruitful field that has only just begun to be explored, but neuroscientists have been calling for this kind of research for decades (Lutz \& Thompson, 2003; Varela, 1996). Since mind-wandering typically consists in a drift of task-focused cognition towards task-unrelated content, it is often accompanied by a drop in performance in the on- going task. Nevertheless, between 30\% and 50\% of waking thoughts are unrelated to ongoing activities (Klinger \& Cox, 1987) and some evidence show that mind- wandering can benefit creativity (Baird et al., 2012). What could be the adaptive value of mind-wandering? Although it's true that mind-wandering takes place during every conceivable activity (Killingsworth \& Gilbert, 2010), the rates of mind-wandering are not identical across all activities. In a typical lab experiment, we deliberately give people a very boring task to do so that we induce high rates of mind-wandering and have some psychological content to study and analyze. So, when we find that people mind-wander a lot during a boring task, or during college lectures for example, we need to remember that in many cases the participant's interest in the task or the K. Fox - Wandering along the spectrum of spontaneous thinking ALIUS Bulletin n2 (2018) aliusresearch.org/bulletin 8 lecture is minimal. For example, in the case of university lectures, perhaps listening to the lecture is relevant to some distant goal of doing well on a final exam and getting a high GPA years from now, but it's easy to see how such distant and relatively vague goals can lose out in the competition with more immediate concerns like interpersonal relationships and conflicts, what to eat, how to schedule your day, and so on. Although I'm not aware of any study that has directly addressed this question, I would say it's a safe bet that motivation and interest in a given task will be strongly inversely correlated with the amount of mind-wandering taking place. A good example is flow states, where people are fully engaged and at an optimal level of difficulty for their given skill level. In the flow state, people hardly report any thinking at all; in fact they often report that any sense of self whatsoever essentially disappears, and they are fully and completely immersed in the activity (Csikszentmihalyi, 2014). But of course much of everyday life isn't like this: we tend to have jobs employing us to do things we are good at, and most things we are good at are by definition overpracticed and hence not fully engaging. And then there are all the daily chores: cooking, cleaning, and so on. These things just need to be done, but they often don't require your full engagement; it's not surprising to me that the brain often disengages from the perceptual world during such activities and instead focuses on other concerns and emotions. I see this as a very clever and resourceful adaptation: if your immediate environment and activities don't require your full attention, your brain very quickly and naturally turns to the recent past or immediate future, or even more distant memories and hypothetical futures, and considers these instead. This seems like it would have clear benefits for everyday life. And if your current goals and concerns happen to be, say, an artistic or scientific problem you are working on, your brain will turn to these concerns-I see the relationship with creativity as just a special case of the more general phenomenon of an automatic, natural focus on current concerns whenever there are cognitive resources available. If you're an artist, your concern will be your next creation; if you're a new parent, your newborn child. In both cases, your mind will tend to focus on your central concerns whenever it has the chance. How is the study of mind-wandering informative regarding clinical conditions like ADHD or for mental well-being in general? Well, if you accept the hypothesis I've advanced above, that spontaneous thought and mind-wandering are a natural, healthy, and probably useful function of the human mind, then you can then conceive numerous mental health conditions as I see the relationship with creativity as just a special case of the more general phenomenon of an automatic, natural focus on current concerns whenever there are cognitive resources available. " " K. Fox - Wandering along the spectrum of spontaneous thinking ALIUS Bulletin n2 (2018) aliusresearch.org/bulletin 9 dysfunctions of this natural propensity toward spontaneous thinking detached from the here and now. Various clinical conditions will distort or exaggerate given aspects of spontaneous thinking, leading to pathological states and experiences that cause significant distress and life disruption for the individual. For instance, we often spontaneously think about the past, and there doesn't seem to be anything wrong with that; it's nice to remember pleasant experiences, and it's probably useful to recall bad experiences in order to process what can be learned from them and how they might be avoided in the future. But in, say, post-traumatic stress disorder, this tendency to spontaneously recall past experiences becomes extremely distressing and overwhelming. It becomes intrusive, in that the memories can't be ignored and instead completely dominate one's present experience and disrupt one's ability to carry on with normal life; and they become repetitive, focusing over and over again on the traumatic experience, rather than canvassing a wide range of times and topics in one's life, as happens with spontaneous memory recall in healthy people. We can look at depressive rumination in a similar way: rather than focusing on a particular traumatic memory, however, one engages repetitively in a series of negative thoughts about the self, one's past failings, and one's dismal prospects for the future. The normal tendency for the mind to spontaneously have thoughts with a wide range of emotional valence that errs on the positive side (Fox et al., under review; Fox et al., 2014) instead becomes intensely focused only on the negative side of things; the normal tendency to think about the self and its relationship to others becomes distorted to emphasize only the worst possible features and outcomes. On the other hand, some people can become "biased" toward emotionally positive and creative thoughts; for instance, there is evidence that practicing meditation and being mindful can nudge spontaneous thought in this direction (Brown \& Ryan, 2003; Fox et al., under review; Frewen et al., 2008; Jazaieri et al., 2015). So on the one hand, I think we need to see overall thought tendencies and content as a very individual, almost trait-like quality, because people differ enormously in what they think about and these differences seem quite stable over time (Fox, 2016); but on the other hand, these baseline tendencies are clearly malleable: they can be skewed toward the negative by various mental health conditions, and perhaps they can be pushed in more positive directions by practices like mindfulness meditation. Of course all these factors are very relevant to general well-being. My colleagues and I (Andrews- Hanna et al., 2013; Christoff et al., 2016), as well as many other researchers, are continuing to work on these problems in an effort to better understand them, but this is a tough undertaking. We still understand very little about what spontaneous thought is, what causes mental illness, and what conduces to general well-being, so understanding how all these things are interrelated is necessarily a long-term project that is still in its infancy. K. Fox - Wandering along the spectrum of spontaneous thinking ALIUS Bulletin n2 (2018) aliusresearch.org/bulletin 10 Because mind-wandering represents a spontaneous cognitive phenomenon unrelated to task-demands, it is challenging to probe and characterize experimentally. What has been the benefit of neuroimaging for the study of mind- wandering and how has it allowed us to derive specific hypotheses about what mind- wandering consists of? Well, the clearest and most foundational finding is that mind-wandering is definitely tied to recruitment of the default network (Fox et al., 2015), which has helped us understand what is going on (psychologically and neurally) in the "resting" state, and forced neuroimagers to reconsider what state they use as a baseline and comparison state for other cognitive tasks. One of the most important insights from neuroimaging, in my view, is that spontaneous thought recruits "executive" brain regions (Christoff et al., 2009; Fox et al., 2015) that are clearly involved in top-down control of attention, metacognition, and so on, and are often thought of as sort of the opposite of what you would expect during mind-wandering. But this finding actually dovetails well with the fact that the wandering mind tends to focus on goals, concerns, future plans, and so on (as we discussed above). This is a nice example of the neuroimaging evidence supporting what has been found from first-person reports and introspective assessment; subjective reports have many limitations and potential biases, so it is encouraging when what people tell us about these private, unverifiable experiences is in fact supported by what's going on in their brains. Another important finding is the similarities between the neural correlates of dreaming and waking mind-wandering: again, first-person reports suggest a lot of similarities, but this doesn't prove that both processes are sharing neural mechanisms. But the neuroimaging evidence, and even evidence from brain lesion patients, strongly suggests that this is indeed the case (Domhoff, 2011; Domhoff \& Fox, 2015; Fox et al., 2016; Fox et al., 2013). As we touched on earlier, such a "continuous" view of spontaneous thought has important implications for understanding different states of consciousness as well as the origins of mental health conditions that involve dysfunctions in spontaneous thinking. So, neuroimaging has already contributed to our understanding in many ways even though there have been very few studies to date. But this is changing rapidly: spontaneous thought is becoming a popular and acceptable topic, and I expect our understanding to advance by leaps and bounds as more, and more sophisticated, neuroimaging studies of spontaneous thought are conducted. We are already starting to go beyond these broad correlations and correspondences and starting to test more specific hypotheses (Kucyi et al., 2013), to investigate the neural correlates of particular types of thought content (Gorgolewski et al., 2014), and to understand the neural basis of how emotions color spontaneous thoughts (Tusche et al., 2014). K. Fox - Wandering along the spectrum of spontaneous thinking ALIUS Bulletin n2 (2018) aliusresearch.org/bulletin 11 Where do you see the study of mind-wandering evolving in the coming years? A major, long-term project will be understanding its relationship(s) to mental illness and general well-being, as we already discussed. A strong personal interest of mine is investigating just how stable people's patterns of thought are, and how these relate to personality and creativity. In contrast, how malleable are spontaneous thought patterns? Can we steer people away from the negative biases that we see in mental illness, and instead nudge them toward positive, constructive, and creative patterns of thinking? K. Fox - Wandering along the spectrum of spontaneous thinking ALIUS Bulletin n2 (2018) aliusresearch.org/bulletin 12 References Andrews-Hanna, J. R., Kaiser, R. H., Turner, A. E. J., Reineberg, A., Godinez, D., Dimidjian, S., \& Banich, M. (2013). A penny for your thoughts: Dimensions of thought content and relationships with individual differences in emotional well- being. Frontiers in Perception Science. Baird, B., Smallwood, J., Mrazek, M. D., Kam, J. W., Franklin, M. S., \& Schooler, J. W. (2012). Inspired by distraction: mind wandering facilitates creative incubation. Psychological Science, 23(10), 1117-1122. Bayne, T., Hohwy, J., \& Owen, A. M. (2016). Are there levels of consciousness?. Trends in cognitive sciences, 20(6), 405-413. Braun, A. R., Balkin, T., Wesensten, N., Carson, R., Varga, M., Baldwin, P., . . . Herscovitch, P. (1997). Regional cerebral blood flow throughout the sleep-wake cycle. Brain, 120(7), 1173-1197. Brown, K. W., \& Ryan, R. M. (2003). The benefits of being present: mindfulness and its role in psychological well-being. Journal of personality and social psychology, 84(4), 822. Buchsbaum, M. S., Gillin, J. C., Wu, J., Hazlett, E., Sicotte, N., Dupont, R. M., \& Bunney, W. E. (1989). Regional cerebral glucose metabolic rate in human sleep assessed by positron emission tomography. Life sciences, 45(15), 1349-1356. Christoff, K., Gordon, A. M.,Smallwood, J., Smith, R., \& Schooler, J. W. (2009). Experience sampling during fMRI reveals default network and executive system contributions to mind wandering. Proceedings of the National Academy of Sciences, 106(21), 8719-8724. Christoff, K., Irving, Z. C., Fox, K. C. R., Spreng, R. N., \& Andrews-Hanna, J. R. (2016). Mind-wandering as spontaneous thought: A dynamic framework. Nature Reviews Neuroscience, 17, 718-731. Csikszentmihalyi, M. (2014). Toward a psychology of optimal experience Flow and the foundations of positive psychology (pp. 209-226): Springer. Domhoff, G. W. (2011). The neural substrate for dreaming: is it a subsystem of the default network? Conscious Cogn, 20(4), 1163-1174. doi: 10.1016/j.concog.2011.03.001 Domhoff, G. W. (2017). The emergence of dreaming: Mind-wandering, embodied simulation, and the default network. Oxford University Press. Domhoff, G. W., \& Fox, K. C. R. (2015). Dreaming and the default network: A review, synthesis, and counterintuitive research proposal. Conscious Cognition, 33, 342-353. Ellamil, M., Fox, K. C. R., Dixon, M. L., Pritchard, S., Todd, R. M., Thompson, E., \& Christoff, K. (2016). Dynamics of neural recruitment surrounding the spontaneous K. Fox - Wandering along the spectrum of spontaneous thinking ALIUS Bulletin n2 (2018) aliusresearch.org/bulletin 13 arising of thoughts in experienced mindfulness practitioners. Neuroimage, 136, 186- 196. Fox, K. C. R. (2016). Functional neuroanatomy of self-generated thought: investigating general brain recruitment, specific neural correlates, and neural origins using functional magnetic resonance imaging and diffusion tensor imaging. (PhD), University of British Columbia, Vancouver, Canada. Fox, K. C. R., Andrews-Hanna, J. R., C., M., Dixon, M. L., Markovic, J., Thompson, E., \& Christoff, K. (under review). Affective neuroscience of undirected thought. Annals of the New York Academy of Sciences, Special Issue: The Year in Cognitive Neuroscience. Fox, K. C. R., Andrews-Hanna, J. R., \& Christoff, K. (2016). The neurobiology of self- generated thought from cells to systems: Integrating evidence from lesion studies, human intracranial electrophysiology, neurochemistry, and neuroendocrinology. Neuroscience, 335, 134-150. Fox, K. C. R., \& Christoff, K. (2014). Metacognitive Facilitation of Spontaneous Thought Processes: When Metacognition Helps the Wandering Mind Find Its Way. The Cognitive Neuroscience of Metacognition (pp. 293-319): Springer. Fox, K. C. R., \& Christoff, K. (2015). Transcranial direct current stimulation to lateral prefrontal cortex could increase meta-awareness of mind wandering. Proceedings of the National Academy of Sciences, 112(19), E2414 Fox, K. C. R., \& Girn, M. (in press). Neural correlates of self-generated imagery and cognition throughout the sleep cycle. In K. C. R. Fox \& K. Christoff (Eds.), The Oxford Handbook of Spontaneous Thought. New York: Oxford University Press. Fox, K. C. R., Nijeboer, S., Solomonova, E., Domhoff, G. W., \& Christoff, K. (2013). Dreaming as mind wandering: Evidence from functional neuroimaging and first- person content reports. Front Hum Neurosci, 7, 412. Fox, K. C. R., Spreng, R. N., Ellamil, M., Andrews-Hanna, J. R., \& Christoff, K. (2015). The wandering brain: Meta-analysis of functional neuroimaging studies of mind- wandering and related spontaneous thought processes. NeuroImage, 111, 611-621. Fox, K. C. R., Thompson, E., Andrews-Hanna, J. R., \& Christoff, K. (2014). Is thinking really aversive? A commentary on Wilson et al.'s "Just think: The challenges of the disengaged mind". Frontiers in Psychology: Cognition. Fox, K. C. R., Zakarauskas, P., Dixon, M. L., Ellamil, M., Thompson, E., \& Christoff, K. (2012). Meditation experience predicts introspective accuracy. PLoS ONE, 7(9), e45370. Frewen, P. A., Evans, E. M., Maraj, N., Dozois, D. J., \& Partridge, K. (2008). Letting go: Mindfulness and negative automatic thinking. Cognitive therapy and research, 32(6), K. Fox - Wandering along the spectrum of spontaneous thinking ALIUS Bulletin n2 (2018) aliusresearch.org/bulletin 14 758-774. Gorgolewski, K. J., Lurie, D., Urchs, S., Kipping, J. A., Craddock, R. C., Milham, M. P., . . . Smallwood, J. (2014). A correspondence between individual differences in the brain's intrinsic functional architecture and the content and form of self-generated thoughts. PLoS ONE, 9(5), e97176. Heiss, W.-D., Pawlik, G., Herholz, K., Wagner, R., \& Wienhard, K. (1985). Regional cerebral glucose metabolism in man during wakefulness, sleep, and dreaming. Brain Research, 327(1), 362-366. Jazaieri, H., Lee, I. A., McGonigal, K., Jinpa, T., Doty, J. R., Gross, J. J., \& Goldin, P. R. (2015). A wandering mind is a less caring mind: Daily experience sampling during compassion meditation training. The Journal of Positive Psychology(ahead-of-print), 1-14. Kam, J. W., Dao, E., Stanciulescu, M., Tildesley, H., \& Handy, T. C. (2013). Mind wandering and the adaptive control of attentional resources. Journal of Cognitive Neuroscience, 25(6), 952-960. Kam, J. W., \& Handy, T. C. (2013). The neurocognitive consequences of the wandering mind: a mechanistic account of sensory-motor decoupling. Frontiers in Psychology, 4. Killingsworth, M. A., \& Gilbert, D. T. (2010). A wandering mind is an unhappy mind. Science, 330, 932. Klinger, E. (2008). Daydreaming and Fantasizing: Thought Flow and Motivation. In K. D. Markman, W. M. P. Klein \& J. A. Suhr (Eds.), Handbook of Imagination and Mental Simulation (pp. 225-239). New York: Psychology Press. Klinger, E. (2013). Goal Commitments and the content of thoughts and dreams: basic principles. Frontiers in Psychology, 4. Klinger, E. \& Cox, W. M. (1987). Dimensions of thought flow in everyday life. Imagination, Cognition and Personality, 7: 105-128 Klinger, E., \& Cox, W. M. (2004). Motivation and the theory of current concerns. Handbook of motivational counseling: Concepts, approaches, and assessment, 3-27. Klinger, E., \& Cox, W. M. (2011). Motivation and the goal theory of current concerns. Handbook of motivational counseling: Goal-based approaches to assessment and intervention with addiction and other problems, 1-47. Kucyi, A., Salomons, T. V., \& Davis, K. D. (2013). Mind wandering away from pain dynamically engages antinociceptive and default mode brain networks. Proceedings of the National Academy of Sciences, 110(46), 18692-18697. K. Fox - Wandering along the spectrum of spontaneous thinking ALIUS Bulletin n2 (2018) aliusresearch.org/bulletin 15 Lutz, A., \& Thompson, E. (2003). Neurophenomenology integrating subjective experience and brain dynamics in the neuroscience of consciousness. Journal of consciousness studies, 10(9-10), 31-52. Madsen, P., Schmidt, J., Wildschiodtz, G., Friberg, L., Holm, S., Vorstrup, S., \& Lassen, N. (1991). Cerebral O2 metabolism and cerebral blood flow in humans during deep and rapid-eye-movement sleep. Journal of Applied Physiology, 70(6), 2597-2601. Maquet, P., Dive, D., Salmon, E., Sadzot, B., Franco, G., Poirrier, R., . . . Franck, G. (1990). Cerebral glucose utilization during sleep-wake cycle in man determined by positron emission tomography and [18 F] 2-fluoro-2-deoxy-d-glucose method. Brain Research, 513(1), 136-143. Nielsen, T. A. (2000). A review of mentation in REM and NREM sleep:"covert" REM sleep as a possible reconciliation of two opposing models. Behavioral and Brain Sciences, 23(06), 851-866. Nir, Y., \& Tononi, G. (2010). Dreaming and the brain: from phenomenology to neurophysiology. Trends in cognitive sciences, 14(2), 88-100. Schooler, J. W., Smallwood, J., Christoff, K., Handy, T. C., Reichle, E. D., \& Sayette, M. A. (2011). Meta-awareness, perceptual decoupling and the wandering mind. Trends in cognitive sciences, 15(7), 319-326. Sze, J. A., Gyurak, A., Yuan, J. W., \& Levenson, R. W. (2010). Coherence between emotional experience and physiology: Does body awareness training have an impact? Emotion, 10(6), 803-814. Tusche, A., Smallwood, J., Bernhardt, B. C., \& Singer, T. (2014). Classifying the wandering mind: Revealing the affective content of thoughts during task-free rest periods. Neuroimage, 97, 107-116. Varela, F. J. (1996). Neurophenomenology: A methodological remedy for the hard problem. Journal of consciousness studies, 3(4), 330-349. Walker, M. P. (2009). The role of sleep in cognition and emotion. Annals of the New York Academy of Sciences, 1156(1), 168-197. Walker, M. P., Liston, C., Hobson, J. A., \& Stickgold, R. (2002). Cognitive flexibility across the sleep-wake cycle: REM-sleep enhancement of anagram problem solving. Cognitive Brain Research, 14(3), 317-324.\par
  \endgroup
}%
\end{document}
